% s-204-p204-m002; original PDF 204
\Aphorism{s-204-p204-m002}{292}
Das grammaticalische Hervortreten des Subjects steht im umgekehrten Verhältniss zu seiner Fülle und Mächtigkeit im Leben: wer Ich thut, masst sich nicht so leicht an, Ich zu sagen. Die Römer berichten von ihren Thaten im Passivum, und bei uns ausgeblasenen Eierschalen von Persönlichkeiten des neunzehnten Jahrhunderts ist das „Ich“ gar ein philosophisches Substantivum geworden! Wir bilden uns zuletzt auf unsere \Foreign{pluralité} \Foreign{du moi} noch etwas ein, obwohl wir sie billig, nämlich um ein blosses Wort- und Begriffsspiel, erstanden haben. Freilich, der moderne Mensch hat an Polypsychie die faustische Zweizahl längst überschritten und beginnt in allem Flor eines siebenfarbigen Regenbogens zu glänzen; ja wenn er alle spectroskopischen Übergänge und Nüancen mitzählen will, mag er sich gleich der Farbenscala eine \Sperr{unendliche} Mannigfaltigkeit dünken — aber wer heisst ihn, jede Spectrallinie als besonderes Ich zu isoliren und zu individualisiren? Wenn die Seele in Rotation geräth und alle ihre Seiten und Winkel von der Umwelt bestrahlen lässt, so steht es \Foreign{verbaliter} frei, diese Seiten- und Winkelseelen als selbständige und autonome Gebilde von ihrer Mutterseele abzulösen; aber sachlich ist damit nichts bewiesen als eben die geschwächte Intensität der modernen Persönlichkeit. Wir haben sieben Achtelseelen, das ist immer noch weniger als eine ganze.
\par

